% Mục 6.1 (bản rút gọn). Rút từ sections/results_lock.tex; nguồn từng số ghi ở đó (results/apple-m5/lock.csv,
% lock_yield.csv, sleep_latency.csv; scripts/report_numbers/rlock_analyze.py, rlock_analyze2.py, extra_analyze.py).
\subsubsection{RQ2: Lock mở rộng thế nào theo số thread, tranh chấp và độ dài critical section?}

\begin{rqanswer}\textbf{Trả lời RQ2.} TTAS mở rộng tốt nhất trong ba lock tự viết. Ở 20 thread
($N = 1009$, $\mathrm{OW} = 0$), nó chậm 2,48 lần cận dưới. MCS gãy mạnh tại số core (13 lần ở 10
thread, 7 lần khi tải nền thấp) và sụp khi vượt số core (1\,700--6\,900 lần). Các ô đó đo với số vòng lặp giảm 100 lần và chỉ cho bậc độ
lớn. Lock thư viện không sụp. \PCname{std_mutex} và \PCname{pthread_mutex} chậm hơn spinlock khi
critical section dài cỡ 0,5~$\mu$s (3,9--4,0 lần cận dưới ở 20 thread, so với 2,48 của TTAS). Riêng
\PCname{os_unfair_lock} giữ 1,6--1,7 lần. Khi critical section chỉ vài nanosecond, cả ba lock thư viện
nhanh hơn spinlock không backoff từ 3 thread. Số liệu của \PCname{apple-m5}. RQ6 kiểm lại trên hai máy còn lại.\end{rqanswer}
% nguồn: rlock_analyze.py, độ chậm N=1009 OW=0: p=20 ttas 2,48, std_mutex 3,93, pthread_mutex 3,97,
% os_unfair_lock 1,59..1,69; mcs p=10 13,0, p=12/16/20 1,70e3/5,16e3/6,88e3. N=11 OW=0 p>=3.

\PCfigure{apple-m5}{lock_slowdown}

\emph{Bằng chứng.} Ở $\mathrm{OW} = 0$, ba spinlock tự viết chậm không quá 1,24 lần cận dưới tới 4
thread (Hình~\ref{fig:apple-m5:lock_slowdown}). Chúng gãy từ 4 lên 6 thread (TTAS 1,08 lên 1,31, MCS
1,24 lên 2,98), tức tại số P-core chứ không tại số core. Từ 8 thread MCS chậm nhất trong mọi lock, chậm
13 lần ở đúng 10 thread rồi sụp hẳn khi vượt số core. TAS và TTAS chỉ chậm dần, tới 3,48 và 2,48 lần ở
20 thread. Lock thư viện đã chậm 2,5--2,9 lần từ 3--4 thread (\PCname{os_unfair_lock} 1,6 lần từ 2
thread), rồi gần như phẳng từ 8 thread (3,5--4,0 lần).
% nguồn: rlock_analyze.py, độ chậm N=1009 OW=0 p=4/6: ttas 1,08/1,31, mcs 1,24/2,98; mcs p=10 13,0;
% p=20 tas 3,48, ttas 2,48; std_mutex p=3/4 2,79/2,93; pthread_mutex 2,45/2,86; p=8..20 3,54..3,97.

Ô $\mathrm{OW} = 15$ là một thí nghiệm đối chứng, vì với $p \le 15$ lock bận tối đa $p/16$ thời gian.
Vậy mà ở 6 thread cả mười hai lock, kể cả lock thư viện, cùng chậm 1,53--2,03 lần cận dưới. Ở
$\mathrm{OW} = 1$ và 2 thread, \PCname{std_mutex} và \PCname{pthread_mutex} đã chậm 2,9 lần cận dưới,
còn ba spinlock tự viết chậm 1,11--1,27 lần, nên lock thư viện bão hoà sớm hơn spinlock. Độ dài critical
section quyết định tỉ lệ giữa chi phí chuyển giao và phần việc. Với $N = 10\,007$ (khoảng 5,2~$\mu$s),
TAS và TTAS gần như trùng nhau. Với $N = 11$ (4,6~ns), từ 3 thread lock thư viện chậm 2,9--6,0 lần cận
dưới, còn TAS chậm 5,2--80,5 lần và TTAS 7,3--28,2 lần.
% nguồn: rlock_analyze.py, độ chậm N=1009 OW=15 p=6: 1,53..2,03; OW=1 p=2: std_mutex 2,86, pthread_mutex
% 2,91, tas/ttas/mcs 1,12/1,11/1,27; OW=0 N=10007 p=20 tas 2,55, ttas 2,53; N=11 OW=0 p>=3: std_mutex
% 4,08..5,96, pthread_mutex 3,68..5,14, os_unfair_lock 2,87..5,02; tas 5,18..80,5, ttas 7,28..28,2.

\emph{Vì sao.} Nhóm cho rằng chỗ gãy ở 4 thread đến từ máy, không từ thuật toán. Ở ô đối chứng, mọi
lock cùng chậm như nhau, và độ công bằng ở 6 thread là 0,84--0,98, tức mọi thread cùng chậm đi. Điều này
phù hợp với việc scheduler luân chuyển thread giữa P-core và E-core chậm hơn, hoặc với tải nền. Trên
\PCname{C4M7Y670NT} (6 P-core), điểm gãy dời sang 6 lên 8 thread (RQ6).

MCS trao lock cho đúng một thread kế tiếp theo thứ tự FIFO. Nếu thread đó không có core, lock nằm yên
cho tới khi scheduler cho nó chạy lại, và cả hàng chờ mất một lượt lập lịch (thời gian scheduler cho
một thread chạy liền trước khi xét lại, cỡ mili giây). TAS và TTAS không bị như vậy, vì thread nào đang chạy cũng lấy được lock rảnh. Chương trình \texttt{handoff\_gap} kiểm cơ chế này.
Nó chạy đúng vòng lặp của phần lock ($N = 1009$, $\mathrm{OW} = 0$) cho MCS, TAS, TTAS và
\PCname{std_mutex} ở 4, 8, 10 và 12 thread, mỗi cấu hình 5 lần. Thread lấy lock đo \emph{khoảng trống
trao lock}, tức thời gian từ lúc lock được nhả tới lúc nó có chủ mới. Ở 4 thread, trung vị của khoảng
trống dài nhất trong một lần chạy là 6--22~$\mu$s với cả bốn lock, nên một khoảng trống dài hơn
100~$\mu$s được coi là lock đã chờ một thread không có core.
% nguồn: experiments/run_extra.py (part_a); handoff.csv; extra_analyze.py (max_gap_us p=4: 6..22).

\begin{table}[H]\centering\small
\begin{tabular}{@{}lcccc@{}}
\toprule
 & 4 thread & 8 thread & 10 thread & 12 thread \\
\midrule
\PCname{mcs}       & 0 / 0    & 0,057 / 8,3 & 0,25 / 47 & 27 / 99 \\
\PCname{tas}       & 0 / 0    & 0 / 0       & 0 / 0     & 0,003 / 0,3 \\
\PCname{ttas}      & 0 / 0    & 0 / 0       & 0 / 0     & 0 / 0 \\
\PCname{std_mutex} & 0 / 0    & 0 / 0       & 0 / 0     & 0 / 0 \\
\bottomrule
\end{tabular}
\caption{Lần trao lock dài hơn 100~$\mu$s trên \PCname{apple-m5}, $N = 1009$, $\mathrm{OW} = 0$: tỉ lệ
trên mọi lần trao lock (\%) / phần thời gian chạy nằm trong các lần đó (\%). Trung vị của 5 lần chạy.}
\label{tab:handoff}
\end{table}
% nguồn: scripts/report_numbers/extra_analyze.py (cột long% và time_in_long%).

Bảng~\ref{tab:handoff} cho thấy MCS ở 10 thread chỉ có 0,25\,\% lần trao lock dài như vậy, nhưng chúng
chiếm 47\,\% thời gian chạy. Ở 12 thread, 27\,\% lần trao lock là lần dài và chúng chiếm 99\,\% thời gian.
Số lần chuyển ngữ cảnh không tự nguyện của tiến trình MCS tăng từ 60 ở 4 thread lên 1\,970 ở 10 thread.
Phép đo này chạy khi load average khoảng 2,8, nên ngay ở 10 thread đã có thread bị tước core. Đối chứng là
TAS ở 10 thread. Tiến trình TAS cũng bị chuyển ngữ cảnh 1\,022 lần, cùng bậc với MCS, nhưng trung vị số
lần trao lock dài bằng 0. Phép đo này không kiểm điểm gãy từ 4 lên 6 thread, nên cơ chế của điểm gãy đó
vẫn là giả thuyết.
% nguồn: extra_analyze.py: mcs p=10 long 0,245 %, time_in_long 47,3 %; p=12 26,8 % / 99,4 %; nivcsw mcs
% p=4/10: 60 / 1\,970; load1 2,8--2,9; tas p=10 nivcsw trung vị 1\,022, long_handovers trung vị 0.

Lock thư viện cho thread chờ ngủ trong kernel, nên thread chờ không giành core của chủ lock và lock không
sụp. Nhóm giả thuyết rằng mỗi lần đánh thức một thread qua kernel tốn một chi phí cố định cỡ 1,5~$\mu$s,
ước lượng từ chính số đo lock. Ở $N = 1009$, mỗi critical section của \PCname{std_mutex} ở 20 thread tốn
thêm 1\,459~ns so với khi không tranh chấp. Ở $N = 10\,007$, mức tăng là 1\,469~ns, gần như không đổi dù
phần việc dài gấp 10 lần. Tài liệu của \PCname{os_unfair_lock} (\texttt{os/lock.h}) ghi rằng thread vừa
nhả có thể lấy lại lock trước khi thread được đánh thức kịp thử. Một critical section của nó chỉ tốn thêm
302--351~ns từ 2 tới 20 thread, và độ công bằng giảm từ 0,97 ở 3 thread xuống 0,20 ở 20 thread. TAS thua
TTAS nhiều nhất khi lock đổi chủ ở mỗi lượt, vì phép nhả lock của TAS phải tranh cache line với các
\texttt{exchange} liên tục của thread chờ. Ở $\mathrm{OW} = 1$ và 8 thread, TAS đạt 0,467 triệu lần lấy
lock mỗi giây, TTAS 1,41 triệu.
% nguồn: rlock_analyze.py, ns/CS std_mutex N=1009 OW=0 p=20: 2\,007; locksingle 548; lock.csv N=10007:
% 6\,847 - 5\,377 = 1\,469. ns/CS os_unfair_lock p=2..20 trừ p=1: 301,7..351,0; fairness p=3 0,97, p=20 0,20.
% Thông lượng OW=1 p=8 ttas 1,41 M/s, tas 0,467 M/s.

\emph{Độ vững.} Cả bảy lần chạy MCS 10 thread đều chậm 7,1--15,5 lần cận dưới. Ở $\mathrm{OW} \ge 1$,
một số ô MCS vượt số core cho kết quả hai cực. Ở 16 thread, $\mathrm{OW} = 15$, bốn lần chạy xong dưới
10~ms, ba lần mất 2,0--10,6~s, nên ở các ô này chỉ khoảng min--max đáng tin.
% nguồn: lock.csv, mcs p=10 OW=0 N=1009: 1,241..2,727 s, cận dưới 0,1755 s; mcs p=16 OW=15: 4 lần < 10 ms,
% 3 lần 2,0..10,6 s.

% ================================================================================================
\subsubsection{RQ3: Backoff giúp hay hại, và lock nhanh nhất có công bằng không?}

\begin{rqanswer}\textbf{Trả lời RQ3.} Khi không có việc ngoài critical section, backoff giúp TAS và gần như không
đổi gì với TTAS. Ở 20 thread
($N = 1009$, $\mathrm{OW} = 0$), backoff bằng \texttt{sleep\_for} đưa TAS từ 12,2~s xuống 3,78~s
(thời gian quy về cùng số vòng lặp). Lợi ích đến từ việc lock gần như không đổi chủ, nên lock nhanh
nhất lại kém công bằng nhất. Khi có việc ngoài critical section và tranh chấp vừa phải, backoff làm
chậm. Hại lớn nhất là ở \PCname{ttas_bo_sleep} khi vượt số core với $\mathrm{OW} = 3$, chậm hơn TTAS
6,1--6,8 lần. Số liệu của \PCname{apple-m5}. RQ6 kiểm lại trên hai máy còn lại.\end{rqanswer}
% nguồn: locktime p=20: tas 12,2; tas_bo_sleep 3,78. rlock_analyze.py ttas_bo_sleep/ttas OW=3 p=16/20: 6,13/6,78.

\PCfigure{apple-m5}{lock_backoff}

\emph{Bằng chứng.} Ở $\mathrm{OW} = 0$ (Hình~\ref{fig:apple-m5:lock_backoff}, cột trái),
\PCname{tas_bo_sleep} và \PCname{tas_bo_yield} chỉ mất 0,31--0,32 lần thời gian của TAS ở 20 thread,
còn \PCname{tas_bo_spin} mất 0,75 lần. Các biến thể của TTAS nằm trong 0,87--1,15 lần TTAS. Ở
$\mathrm{OW} = 3$ và 4--6 thread, backoff làm TAS chậm hơn (\PCname{tas_bo_sleep} 1,32--1,33 lần). Ở 20
thread, \PCname{tas_bo_sleep} và \PCname{tas_bo_yield} lại nhanh hơn (0,69 và 0,55 lần TAS), còn
\PCname{tas_bo_spin} vẫn chậm hơn (1,18 lần).
% nguồn: locktime p=20: 3,78/12,2 = 0,31; 3,95/12,2 = 0,32. rlock_analyze.py tỉ số so với tas N=1009: OW=0 p=20
% tas_bo_spin 0,745; ttas_bo_*/ttas 0,87..1,15. OW=3: tas_bo_sleep p=4/6 1,32/1,33; p=20 0,688, 0,549, 1,18.

\PCfigure{apple-m5}{lock_fairness}

Ở $\mathrm{OW} = 0$, MCS giữ độ công bằng 1 tới 10 thread (Hình~\ref{fig:apple-m5:lock_fairness}). Độ
công bằng của \PCname{tas_bo_sleep} bằng đúng $1/p$, tới 0,05 ở 20 thread, tức là các thread chạy xong
lần lượt từng thread một. TAS và TTAS cũng thiên vị thread vừa nhả lock, TTAS nhiều hơn (0,26 so với
0,49 ở 4 thread). \PCname{os_unfair_lock}, lock thư viện nhanh nhất, kém công bằng nhất trong ba lock thư
viện (0,20 so với 0,78--0,79 ở 20 thread). Khi $\mathrm{OW} > 0$, thread vừa nhả đi gọi
\texttt{is\_prime} thay vì lấy lại lock ngay, nên \PCname{os_unfair_lock} gần công bằng ở
$\mathrm{OW} = 3$ và 15 (0,87--1,00).
% nguồn: rlock_analyze.py, fairness N=1009 OW=0: mcs p=2..10 1,00; tas_bo_sleep p=20 0,05; tas/ttas p=4 0,49/0,26;
% os_unfair_lock p=20 0,20; std_mutex 0,78, pthread_mutex 0,79. lock.csv os_unfair_lock OW=3, 15: 0,87..1,00.

\emph{Vì sao.} Các biến thể \PCname{tas_bo_*} lùi lại sau mọi lần \texttt{exchange} thất bại, kể cả khi
lock chỉ đang bị giữ, nên thread chờ gần như luôn ngủ hoặc nhường core. Các biến thể \PCname{ttas_bo_*}
chỉ lùi lại sau một lần \texttt{exchange} thua và vẫn đọc liên tục như TTAS. Một lần \texttt{sleep\_for}
yêu cầu 100~ns thật ra ngủ trung vị 2\,670~ns. Ở $\mathrm{OW} = 0$, chủ lock nhả rồi lấy lại lock ngay ở
lượt kế tiếp. Thread vừa thức dậy gần như luôn thấy lock bận và ngủ tiếp với thời gian yêu cầu gấp đôi,
tới khoảng 100~$\mu$s. Vì vậy một thread giữ lock tới hết vòng lặp của nó, rồi mới tới thread khác, đúng
như độ công bằng $1/p$. \PCname{tas_bo_spin} chỉ giãn các lần thử, thread chờ vẫn giữ core, nên lợi ích
nhỏ hơn. Khi $\mathrm{OW} > 0$, lock rảnh trong lúc thread chờ còn ngủ ít nhất 2,5~$\mu$s, và ở tranh
chấp vừa phải đó là thời gian mất. Khi vượt số core, thread ngủ hoặc nhường core để lại core cho thread
đang làm việc ngoài critical section. Nhóm cho rằng đó là lý do \PCname{tas_bo_sleep} và
\PCname{tas_bo_yield} nhanh hơn TAS ở 20 thread.
% nguồn: sleep_latency.csv: sleep_for 100 ns trung vị 2\,667 ns, nhỏ nhất 2\,541 ns; include/backoff.hpp.

Khi vượt số core với $\mathrm{OW} = 3$, \PCname{ttas_bo_sleep} chậm hơn TTAS 6,1--6,8 lần. Nhóm giả
thuyết rằng ở biến thể này một phần thread chờ quay đọc và chiếm core, phần khác ngủ, và thread vừa thức
dậy có thể giành core của chủ lock. \PCname{ttas_bo_yield}, chỉ khác ở chỗ nhường core thay vì ngủ, không
chậm hơn TTAS (1,97~s so với 2,42~s ở 20 thread, thời gian quy về cùng số vòng lặp). Số đo này phù hợp
với giả thuyết.
% nguồn: tables.tex locktime-ow3 p=20: ttas_bo_yield 1,97, ttas 2,42.

\emph{Độ vững.} Kết quả của \PCname{ttas_bo_sleep} khi vượt số core phân tán mạnh. Ở 16 thread và
$\mathrm{OW} = 3$, thời gian quy về cùng số vòng lặp đi từ 1,0 tới 18,0~s, với trung vị 10,4~s. Một lượt phụ chạy lại TAS và
TTAS với lệnh máy \texttt{yield} thay \texttt{isb} trong vòng spin ở 1, 4, 10, 20 thread, và khác biệt không vượt độ phân
tán.
% nguồn: lock.csv ttas_bo_sleep p=16 OW=3: 1,04..18,04 s, trung vị 10,4 s; lock_yield.csv vs lock.csv.
